\subsection{Literatura Vigente}

A literatura sobre vulnerabilidades e priorização de correções converge em três eixos principais: predição de exploração, avaliação crítica do CVSS e integração/qualidade de dados.

Bozorgi et al. (2010) estão entre os trabalhos iniciais a demonstrar que atributos públicos de vulnerabilidades podem ser usados para classificar vulnerabilidades e prever desfechos relacionados à exploração com desempenho superior a heurísticas simples \cite{bozorgi2010_beyond_heuristics}. O trabalho reforça a importância de dados estruturados, históricos e reutilizáveis. Entretanto, seu contexto é anterior ao amadurecimento de abordagens posteriores centradas em exploração observada em ambiente real e em scores probabilísticos padronizados.

Zhang et al. (2011) analisam o uso da NVD para prever vulnerabilidades em software, mais especificamente o tempo até a próxima vulnerabilidade em uma aplicação, e indicam que a base, isoladamente, possui poder preditivo limitado em muitos cenários \cite{zhang2011_nvd_predict,zhang2015_nvd_cyber_risks}. Esse resultado é importante para a presente proposta porque mostra que a NVD, embora central, não deve ser a única fonte de dados para análises de segurança.

Allodi e Massacci (2012; 2014) investigam a relação entre scores de severidade e exploração observada, argumentando que o uso isolado de CVSS como critério de priorização pode produzir baixa efetividade operacional \cite{allodi2012_preliminary,allodi2014_case_control}. Esses estudos são fundamentais porque sustentam a necessidade de enriquecer severidade técnica com sinais de exploração, como PoCs, exploração observada e outros indicadores de ameaça.

Sabottke et al. (2015) exploram sinais públicos, incluindo mídias sociais e metadados de vulnerabilidades, para prever exploração real \cite{sabottke2015_social_media}. Embora o uso de dados sociais apresente desafios de reprodutibilidade e volatilidade, o trabalho reforça a hipótese de que sinais externos aos registros técnicos tradicionais podem melhorar a priorização.

Jacobs et al. introduzem a linha de trabalho de predição de exploração para apoio à remediação em artigo de 2020 e formalizam o EPSS em artigo de 2021, após uma proposta inicial divulgada em 2019 \cite{jacobs2020_remediation,jacobs2019_epss,jacobs2021_epss}. Essa contribuição representa uma mudança importante: em vez de priorizar apenas por severidade, passa-se a considerar a probabilidade de uma vulnerabilidade ser explorada. Para o dataset proposto, o EPSS é uma camada essencial de enriquecimento, ainda que deva ser interpretado como componente de ameaça, não como medida completa de risco.

Jiang et al. (2021) analisam inconsistências em repositórios de vulnerabilidades e mostram que problemas de qualidade de dados podem redirecionar incorretamente esforços de segurança \cite{jiang2021_inconsistency}. Essa conclusão se relaciona diretamente com a proposta deste trabalho, que preserva proveniência e busca preparar a reconciliação entre fontes heterogêneas.

Trabalhos recentes sobre knowledge graphs e relações entre CVE, CWE e CPE reforçam a importância de integrar identificadores e modelar relações entre vulnerabilidades, fraquezas e produtos \cite{shi2024_threat_kg,yin2024_vulkg}. Esses estudos indicam que a integração de bases não é apenas uma tarefa de engenharia, mas uma condição para permitir análises mais robustas, auditáveis e reprodutíveis.

De forma complementar, estudos recentes \textit{package-centric} e baseados em OSV reforçam a importância de representar vulnerabilidades no nível de ecossistemas, pacotes e versões concretas, além de explicitar traços de propagação e reprodutibilidade do dataset \cite{wu2024_affected_libraries,ruohonen2025_osv_propagation,mandl2026_osv_groundtruth}. Isso é especialmente relevante para softwares open source e para estudos que dependem de curadoria precisa de versões vulneráveis e corrigidas.

\begin{table}[htbp]
    \centering
    \caption{Trabalhos relacionados e lacunas exploradas pelo dataset proposto}
    \label{tab:trabalhos-relacionados-lacunas}
    \small
    \renewcommand{\arraystretch}{1.15}
    \resizebox{\textwidth}{!}{%
        \begin{tabular}{|p{2.9cm}|p{2.3cm}|p{3.0cm}|p{3.0cm}|p{3.1cm}|p{3.4cm}|}
            \hline
            \textbf{Referência}                                              & \textbf{Objetivo}                                              & \textbf{Fontes/abordagem}                     & \textbf{Achado principal}                                                        & \textbf{Relação com este trabalho}               & \textbf{Lacuna explorada pelo dataset proposto}                                           \\
            \hline
            Bozorgi et al. (2010)                                            & Prever desfechos de exploração                                 & Atributos públicos e aprendizado de máquina   & Dados públicos podem superar heurísticas simples                                 & Justifica estruturar dados públicos para análise & Contexto anterior a linhas mais recentes de exploração observada e scores probabilísticos \\
            \hline
            Zhang et al. (2011; 2015)                                        & Avaliar NVD para predição de vulnerabilidades em software      & NVD e mineração de dados                      & NVD isolada tem poder preditivo limitado                                         & Reforça necessidade de multi-fonte               & Falta integração com advisories, exploração e scores probabilísticos                      \\
            \hline
            Allodi e Massacci (2012; 2014)                                   & Comparar severidade e exploração observada                     & NVD, Exploit-DB e dados de ataques/exploração & CVSS isolado é fraco como proxy operacional                                      & Justifica combinar CVSS com sinais de exploração & Parte dos sinais usados não é plenamente aberta/reprodutível                              \\
            \hline
            Sabottke et al. (2015)                                           & Prever exploração com sinais públicos                          & Twitter, NVD e ground truth de exploração     & Sinais externos melhoram predição                                                & Reforça valor de variáveis contextuais           & Dados sociais são voláteis e difíceis de reconstruir                                      \\
            \hline
            Jacobs et al. (2019; 2021)                                       & Propor e formalizar EPSS                                       & Modelagem probabilística de exploração        & Probabilidade de exploração complementa severidade                               & Motiva integração de EPSS                        & EPSS mede ameaça, não risco total                                                         \\
            \hline
            Jiang et al. (2021)                                              & Avaliar inconsistência de bases                                & CVE, NVD e fornecedores                       & Inconsistências afetam priorização                                               & Justifica proveniência e reconciliação           & Escopo limitado a atributos e casos observados                                            \\
            \hline
            Shi et al. (2024); Yin et al. (2024)                             & Integrar relações CVE-CWE-CPE e entidades correlatas           & Knowledge graphs                              & Relações integradas revelam associações ausentes e apoiam análise de risco       & Aponta evolução estrutural futura do dataset     & Nem sempre inclui sinais operacionais como KEV/EPSS e advisories package-centric          \\
            \hline
            Wu et al. (2024); Ruohonen e Ramadan (2025); Mandl et al. (2026) & Construir análise package-centric e datasets multi-ecossistema & Bases open source, OSV e curadoria de versões & Ecossistemas, pacotes e versões exigem curadoria precisa e snapshot reprodutível & Motiva integração com GHSA/OSV e versionamento   & Escopo frequentemente centrado em open source e, em parte, ainda emergente                \\
            \hline
        \end{tabular}%
    }
\end{table}
